\documentclass[11pt,a4paper]{amsart}

\usepackage[T1]{fontenc}
\usepackage{lmodern}
\usepackage[english]{babel}

\usepackage{amsmath, amssymb, amsthm}
\usepackage{hyperref}
\hypersetup{hidelinks}
\usepackage{enumerate}
\usepackage{geometry}
\usepackage{float}
\usepackage{csquotes}

% --- Theorem environments ---
\theoremstyle{plain}
\newtheorem{theorem}{Theorem}[section]
\newtheorem{lemma}[theorem]{Lemma}
\newtheorem{corollary}[theorem]{Corollary}
\newtheorem{proposition}[theorem]{Proposition}

\theoremstyle{definition}
\newtheorem{definition}[theorem]{Definition}

\theoremstyle{plain}
\newtheorem{conjecture}[theorem]{Conjecture}

\theoremstyle{remark}
\newtheorem{remark}[theorem]{Remark}
\newtheorem{claim}[theorem]{Claim}

\usepackage[
    backend=biber,
    style=numeric,
    sorting=none
]{biblatex}
\addbibresource{bibliography.bib}

\geometry{margin=2.8cm}
\emergencystretch=3em

% --- Formalization cross-reference macros ---
\newcommand{\lean}[1]{\begingroup\upshape\ttfamily\def\_{\textunderscore\allowbreak}#1\endgroup}
\newcommand{\Formal}[1]{\par\smallskip\noindent\begingroup\raggedright\emph{Formalization.} #1\par\endgroup\smallskip}

\title{The Structural Sieve}
\author{Artur Flamandzki}
\thanks{Email address: \texttt{flamandzki.artur@gmail.com}}
\date{\today}

\begin{document}

\begin{abstract}
We take the prime base, not the interval, as the object of study, and measure the
deterministic reach of its generative sieve. The Bertrand--Chebyshev theorem is then not
a goal but a consequence: asking how far the base reaches collapses the postulate to a
single positivity statement about an exact window count --- the window itself derived from
the base --- which we close in the natural generality of the construction.

Starting from a complete prime base $\mathcal{B}=\{q\ \text{prime}: q<P_{\max}\}$ we
separate two determinisms of the generative sieve. The base is \emph{self-contained}
exactly up to $P_{\min}\cdot P_{\max}$, where $P_{\min}$ is the smallest prime that
participates in the sieve: every composite in the window $(P_{\max}, P_{\min}P_{\max}]$ is
built entirely from the base, and immediately beyond the edge the first composite requiring
a prime \emph{outside} the base appears, namely $P_{\min}\cdot\operatorname{nextprime}(P_{\max})$.
Its \emph{generative decision} --- a survivor is void $\Leftrightarrow$ prime --- reaches
further, throughout the zone $(P_{\max}, P_{\max}^2)$, and only there does it break. Inside the
self-contained window this decision forces every survivor to be a new prime, so the base
cannot cover its own window and a prime outside it is necessary: this forcing is the
\emph{atom}. The self-contained edge is thus \emph{derived} from the base, not assumed. This yields a generalized window
family $(P_{\max}, P_{\min}P_{\max}]$ obtained by \emph{truncated arithmetic}
(sieving with a reduced base); the classical Bertrand window $(n, 2n]$ is the
\emph{tightest} member, $P_{\min}=2$, and every wider member follows from it by
monotonicity. The whole construction is invariant under dilation of scale: the pair
(window, base) is a single self-similar object, and this self-similarity is what the
window count inherits.

The reduction sends non-emptiness of the window to a single positivity statement (the
\emph{atom}), which the proof settles in two regimes of different character. In the
\emph{constructive} regime the mechanism is generative: at every concrete scale it decides the
window outright --- the sieve is finite and its evaluation decidable, so no counterexample can
arise at any fixed step --- and there the theorem is closed by structure alone. What a
generative process cannot produce is a single finite object for the completed infinite, the
uniform statement over all scales at once; that statement is settled not by construction but
existentially. It is exactly here that the central binomial coefficient $\binom{2n}{n}$
supplies closure. It enters not as an external estimate but as the object whose prime
factorization records the window's new sources: a source $q\in(n,2n]$ divides $\binom{2n}{n}$
exactly once. It is present not for any fixed window but because the proof must quantify over
all windows at once, and it is reached from the paper's own purity law, not borrowed. The
extremal window ($P_{\min}=2$) is then closed by comparing the growth bound
$4^{n}\le (2n+1)\binom{2n}{n}$ with the old-source upper bound; both bounds are reproved within
the development, so its main theorem imports no Mathlib Bertrand theorem. These are the same
two inequalities Erd\H{o}s used in 1932: reaching the same object by a different route is a
convergence, with historical priority but not exclusivity. The entire development is
machine-checked in Lean~4 / Mathlib with no \texttt{sorry} and no appeal to Mathlib's proof of
Bertrand's postulate.
\end{abstract}

\maketitle

%======================================================================
\section{Introduction}

Bertrand's postulate asserts that for every integer $n>1$ there is a prime $p$ with
$n<p\le 2n$. Bertrand conjectured it in 1845 on numerical evidence \cite{bertrand1845};
Chebyshev proved it in 1852 by analytic bounds on $\vartheta$ and $\psi$
\cite{chebyshev1852}; Ramanujan gave a short $\Gamma$-function argument in 1919
\cite{ramanujan1919}; and Erd\H{o}s, in 1932, gave the elementary proof that has since
become canonical, extracting the statement from divisibility properties of the central
binomial coefficient $\binom{2n}{n}$ \cite{erdos1932}. Nair later observed that the same
bounds follow from $\operatorname{lcm}(1,\dots,2n)$ \cite{nair1982}. Each of these proofs
stands on its predecessors; none discards the tools it inherits.

The present work stands on the same lineage, and is explicit about where. Its
contribution is not a new inequality but a new \emph{architecture}: we show that
Bertrand's postulate is the tightest instance of a self-similar family of window
statements, that the window and its right endpoint are \emph{forced by the prime base}
rather than posited, and that the entire theorem localizes to a single positivity
statement about an exact count. What remains at that one point is quantitative, and there
we insert the lightest classical certificate --- which turns out to be Erd\H{o}s's pair of
inequalities.

To the author's knowledge, following a literature comparison against the closest related
material --- Jacobsthal's function and its primorial-restricted gap bounds
\cite{jacobsthal1961,iwaniec1978}, and the generalized Bertrand-type window results of
\cite{daspaul2019} --- three things distinguish this construction; \S\ref{sec:provenance}
records the comparison this claim rests on and its limits. First, the window $(P_{\max},2P_{\max}]$ and, more generally,
$(P_{\max},P_{\min}P_{\max}]$, together with the right endpoint $P_{\min}P_{\max}$, are
\emph{derived} from the base $\mathcal{B}=\{q\ \text{prime}: q<P_{\max}\}$: the endpoint
is exactly the buildability edge of the base, and self-containment provably fails one step
beyond it, while the generative decision (void $\Leftrightarrow$ prime) reaches to the square
of the anchor (\S\ref{sec:reach}, \S\ref{sec:edge}). Second, the statements assemble into a
generalized family indexed by the smallest participating prime $P_{\min}$, in which
Bertrand is the tightest ($P_{\min}=2$) and every wider member is a corollary by
monotonicity (\S\ref{sec:family}); the family is obtained by \emph{truncated arithmetic},
i.e.\ sieving with a reduced base, and is manifestly invariant under dilation of scale
(\S\ref{sec:trunc}). Third, non-emptiness of the window collapses to a single exact-count
statement (the atom), and entire regimes of it close with no counting at all
(\S\ref{sec:atom}); the residual case is closed on the central binomial coefficient, whose
prime factorization records the window's content (\S\ref{sec:newton}, \S\ref{sec:erdos}).

We are careful about provenance. The development is neither independent of Erd\H{o}s nor a
disguised copy of him. The structural reduction is independent and shows \emph{why} the
postulate collapses to one atom; the object on which the atom is finally settled --- the
central binomial --- is reached from the internal purity law of the construction, and
would be reached here regardless of the historical record. That Erd\H{o}s settled the same
object in 1932 is a convergence on one mathematical object; he holds the historical
priority of use, but no exclusivity over it. Every claim below is machine-checked in
Lean~4 with Mathlib; the closing certificate is the sole place a Mathlib import about
$\binom{2n}{n}$ is used, and Mathlib's own proof of Bertrand's postulate is never
invoked.

\paragraph{Origin of the construction.}
The construction did not begin from Bertrand's postulate, and it is worth recording the
observation from which it did begin, since that observation fixes the natural home of the
statement. Consider the prime-factor incidence of the integers: for each $n$, record which
primes divide it (Table~\ref{tab:incidence}). Read by rows, the primes are exactly the
positions carrying no mark --- the survivors of the factorization. Read by columns the table
is more telling: a prime $p$ makes its first appearance as a factor at $2p$, its first even
multiple, and only from there does it begin to eliminate positions. Consequently the base of
primes below a scale $P$ already accounts for every composite up to $2P$: inside $(P,2P]$
every eliminated position is struck by a prime $<P$, so every survivor there is a new prime.
What the observation measures is thus the \emph{reach of deterministic factorization} --- how
far a fixed base generates the integers before a genuinely new prime is required --- and that
reach is exactly $2P$. Whether the interval $(P,2P]$ happens to contain a prime was not the
starting point; it is a corollary of the reach, read off afterwards.

\begin{table}[htbp]
\centering
\small
\begin{tabular}{r|ccccc|cccc}
$n$ & 2 & 3 & 5 & 7 & 11 & 13 & 17 & 19 & 23\\ \hline
2  &            &            &            &            &            & & & &\\
3  &            &            &            &            &            & & & &\\
4  & $\bullet$  &            &            &            &            & & & &\\
5  &            &            &            &            &            & & & &\\
6  & $\bullet$  & $\bullet$  &            &            &            & & & &\\
7  &            &            &            &            &            & & & &\\
8  & $\bullet$  &            &            &            &            & & & &\\
9  &            & $\bullet$  &            &            &            & & & &\\
10 & $\bullet$  &            & $\bullet$  &            &            & & & &\\
11 &            &            &            &            &            & & & &\\
12 & $\bullet$  & $\bullet$  &            &            &            & & & &\\
13 &            &            &            &            &            & & & &\\
14 & $\bullet$  &            &            & $\bullet$  &            & & & &\\
15 &            & $\bullet$  & $\bullet$  &            &            & & & &\\
16 & $\bullet$  &            &            &            &            & & & &\\
17 &            &            &            &            &            & & & &\\
18 & $\bullet$  & $\bullet$  &            &            &            & & & &\\
19 &            &            &            &            &            & & & &\\
20 & $\bullet$  &            & $\bullet$  &            &            & & & &\\
21 &            & $\bullet$  &            & $\bullet$  &            & & & &\\
22 & $\bullet$  &            &            &            & $\bullet$  & & & &\\
23 &            &            &            &            &            & & & &\\
24 & $\bullet$  & $\bullet$  &            &            &            & & & &\\
\end{tabular}
\caption{Prime-factor incidence for $2\le n\le 24$. A mark in column $p$ records $p\mid n$;
the empty rows ($n=2,3,5,7,11,13,17,19,23$) are the primes, the survivors of the
factorization. Each prime $p$ first appears as a factor at $2p$, so the primes $\ge 13$ carry
no mark below $2\cdot13=26$: the rule at $P_{\max}=11$ separates the emerged base from the
primes not yet required. The base of primes $<P$ thus covers every composite in $(P,2P]$, and
the reach of the fixed base is exactly $2P$ --- the endpoint derived in \S\ref{sec:edge}.}
\label{tab:incidence}
\end{table}

\paragraph{Independence and naming.}
The right endpoint $2P$, the notion of a complete base, and the generative sieve built on
them (\S\ref{sec:reach}--\S\ref{sec:atom}) were arrived at from this observation,
independently of the classical postulate and of Erd\H{o}s's proof; the coincidence of the
tightest case with the interval of Bertrand's postulate was recognized only afterwards. We
keep the classical name by historical priority, while noting that the statement here is
naturally one about \emph{primes} rather than natural numbers: the object is the window
$(P_{\max},P_{\min}P_{\max}]$ anchored at a prime scale, not $(n,2n]$ anchored at an
arbitrary integer, and the two coincide precisely at the tightest member $P_{\min}=2$
(Corollary~\ref{cor:classical}).

Throughout, $P$ denotes a prime playing the role of the current scale (the anchor
$P_{\max}$); $\mathcal{B}(P)=\{q\ \text{prime}: q<P\}$ is the base below it; and, for a
finite set $S$ of primes, a position $n$ is \emph{$S$-void} if no element of $S$ divides
$n$. All Lean identifiers are typeset as \lean{like\_this} and refer to the accompanying
formalization,\footnote{Publicly available at
\url{https://github.com/Flaman55/StructuralSieve}.}
in which each may be located by name.

%======================================================================
\section{The prime base and its deterministic reach}
\label{sec:reach}

\begin{definition}[Base and void]
For a prime $P$ let the \emph{complete base below $P$} be
$\mathcal{B}(P):=\{q\ \text{prime}: q<P\}$, and let the \emph{Bertrand window} be
$W(P):=(P,2P]\cap\mathbb{Z}$. A position $n$ is \emph{$\mathcal{B}(P)$-void} if
$q\nmid n$ for every $q\in\mathcal{B}(P)$, equivalently if
$\gcd\!\big(n,\prod_{q<P}q\big)=1$.
\end{definition}

\Formal{\lean{isVoid}, \lean{isVoid\_iff\_coprime}, \lean{gps\_window};
$\prod_{q<P}q$ is \lean{primorial\_below}.}

\begin{remark}[One scale parameter fixes base, window, and width]\label{rem:triple}
A single prime $P$ determines all three data: it is the largest prime of the base (the complete
set of primes up to $P$, with none omitted), the left endpoint of the window $W(P)=(P,2P]$, and
its width $|W(P)|=P$. The base and the window are thus both functions of $P$, not independent
parameters. (Whether $P$ itself is counted in the base is immaterial to the covering: its only
multiple in the window is $2P$, already divisible by $2$; see
\lean{anchor\_idle\_in\_own\_window}.) The same structure underlies Bertrand's classical window
$(n,2n]$, with $n$ in the role of $P$: the present statement is phrased on a prime scale rather
than an arbitrary integer, and the two coincide at the extremal member $P_{\min}=2$
(\S\ref{sec:main}). The quantity that governs the statement is accordingly the local deficiency
of the base on its own window --- the length of the longest interval of covered positions in
$W(P)$ --- which external gap functions only bound from above.
\end{remark}

The one arithmetic fact that governs the window is the least-prime-factor bound: a
composite that is not too large must reveal a small factor.

\begin{lemma}[Least prime factor in the window]\label{lem:lpf}
Let $2<P$ and $n\in W(P)$. If $n$ is composite then its least prime factor satisfies
$\operatorname{minFac}(n)\le\sqrt{2P}<P$; in particular
$\operatorname{minFac}(n)\in\mathcal{B}(P)$.
\end{lemma}

\begin{proof}
If $n=ab$ with $1<a\le b$ then $a\le\sqrt n$. For $n\in W(P)$ we have $n\le 2P$, so the
least prime factor is $\le\sqrt{2P}$, and $\sqrt{2P}<P$ for $P>2$.
\end{proof}

\Formal{\lean{least\_prime\_factor\_bound} (\lean{LPF.lean}),
\lean{window\_composite\_minFac\_sq\_le} (\lean{Rings.lean}).}

Lemma~\ref{lem:lpf} turns void into primality inside the window: if a position survives
the base, it has no factor below $P$ and is not composite, hence prime.

\begin{theorem}[Determinism of the window]\label{thm:det}
Let $P$ be prime with $2<P$ and $n\in W(P)$. Then $n$ is $\mathcal{B}(P)$-void if and only
if $n$ is prime.
\end{theorem}

\begin{proof}
If $n$ is prime and $n>P$ then no $q<P$ divides $n$, so $n$ is void. Conversely, a void
$n$ has no prime factor below $P$; were it composite, Lemma~\ref{lem:lpf} would supply
one. Hence $n$ is prime.
\end{proof}

\Formal{\lean{void\_in\_window\_prime}, \lean{void\_iff\_prime\_in\_window},
\lean{void\_iff\_prime\_in\_deterministic\_zone}; the underlying
\enquote{composite $\Rightarrow$ covered by an earlier prime} step is
\lean{composites\_covered\_by\_prev} and \lean{prime\_iff\_uncovered\_by\_prev}
(\lean{ZeroForce.lean}).}

Two consequences frame the rest of the paper. First, non-emptiness of $W(P)$ as a set of
primes is \emph{equivalent} to the existence of a single $\mathcal{B}(P)$-void position;
this is the object we will count. Second, the mechanism of Theorem~\ref{thm:det} does not
depend on $P$: it has the same form at every anchor.

\begin{proposition}[Scale invariance of the mechanism]\label{prop:scale}
For every prime $P$ with $2<P$ and every $n\in W(P)$, the equivalence
$\big(n\ \text{is}\ \mathcal{B}(P)\text{-void}\big)\Leftrightarrow\big(n\ \text{is prime}\big)$
holds verbatim. The rule is invariant under the choice of anchor.
\end{proposition}

\Formal{\lean{mechanism\_scale\_invariant} (\lean{Rings.lean}).}

Proposition~\ref{prop:scale} is the precise content of the informal statement that the
construction \enquote{behaves the same at the start of the number line and at any later
point.} We return to it in \S\ref{sec:family}: the pair $\big(W(P),\mathcal{B}(P)\big)$ is
a single object viewed at scale $P$, and moving along $\mathbb{N}$ rescales window and base
together.

%======================================================================
\section{The buildability edge: why $P_{\min}P_{\max}$}
\label{sec:edge}

The window's right endpoint is not chosen; it is the exact reach of the base. Fix the
smallest prime allowed to participate in the sieve, $P_{\min}$ (for Bertrand,
$P_{\min}=2$), and a scale $P_{\max}$.

\begin{theorem}[Buildability]\label{thm:build}
Let $2\le P_{\min}$ and let $n$ be composite with $\operatorname{minFac}(n)\ge P_{\min}$
and $n\le P_{\min}P_{\max}$. Then every prime factor of $n$ is $\le P_{\max}$; that is, $n$
is fully built from primes in $[P_{\min},P_{\max}]$.
\end{theorem}

\begin{proof}
Let $r$ be any prime factor of $n$ and write $n=r\,s$. Since $n$ is composite,
$s\ge\operatorname{minFac}(n)\ge P_{\min}$, hence
$r=n/s\le P_{\min}P_{\max}/P_{\min}=P_{\max}$.
\end{proof}

\Formal{\lean{rough\_composite\_smooth} (general $P_{\min}$),
\lean{window\_composite\_smooth} (case $P_{\min}=2$).}

Thus below $P_{\min}P_{\max}$ every composite with least factor $\ge P_{\min}$ is
$P_{\max}$-smooth: the base can build it, and a surviving position can only be a new
prime. The edge is sharp.

\begin{proposition}[First unbuildable composite]\label{prop:first}
The smallest composite with least factor $\ge P_{\min}$ that is \emph{not}
$P_{\max}$-smooth is $P_{\min}\cdot\operatorname{nextprime}(P_{\max})$, which exceeds
$P_{\min}P_{\max}$. Hence the reach $P_{\min}P_{\max}$ is exact.
\end{proposition}

For $P_{\min}=2$ the self-contained reach is $2P_{\max}$. Two distinct boundaries meet the base
here and must not be conflated. The base's \emph{self-containment} ends at $2P$: just beyond it
a composite requiring a prime outside the base appears (Proposition~\ref{prop:first}), and $2P$
is the largest interval above $P$ in which no element properly divides another
(Theorem~\ref{thm:self}); this is what fixes the window. The sieve's \emph{generative decision}
--- the equivalence void $\Leftrightarrow$ prime of Theorem~\ref{thm:det} --- reaches much
further, throughout the zone $(P,P^2)$, and breaks only at the square of the anchor.

\begin{theorem}[Generative decision reaches the square scale]\label{thm:break}
Let $P$ be prime, $2<P$. Every $\mathcal{B}(P)$-void $n$ with $P<n<P^2$ is prime, and this
fails at $n=P^2$: the square $P^2$ is composite yet $\mathcal{B}(P)$-void. Hence the generative
decision of Theorem~\ref{thm:det} reaches exactly the zone $(P,P^2)$ --- the square of the
anchor --- far beyond the self-contained edge $2P$.
\end{theorem}

\begin{proof}
If $P<n<P^2$ is $\mathcal{B}(P)$-void then no prime $<P$ divides it, so every prime factor of
$n$ is $\ge P$; were $n$ composite it would be a product of two such factors, hence $\ge P^2$,
against $n<P^2$. So $n$ is prime. For the failure, $P^2$ is composite and its only prime factor
is $P\notin\mathcal{B}(P)$, so it is $\mathcal{B}(P)$-void; the least prime $q>P$ gives a
further composite void $q^2>P^2$.
\end{proof}

\Formal{Positive zone \lean{void\_iff\_prime\_in\_deterministic\_zone}; the break is witnessed
in \lean{determinism\_breaks\_above} (\lean{Rings.lean}), whose exhibited composite void is
$q^2$.}

It is self-containment, not the generative decision, that isolates \emph{why} $2P$ is the
window's true boundary: $(P,2P]$ is exactly the largest interval above $P$ in which no element
properly divides another, so elimination can proceed only through the base and never through a
survivor.

\begin{theorem}[Self-containment, sharp at $2P$]\label{thm:self}
For $m,n\in W(P)$, $m\mid n$ implies $m=n$: no proper divisibility holds inside $W(P)$.
This fails immediately above the edge: $P+1\in W(P)$ but its proper multiple
$2(P+1)>2P$ leaves the window.
\end{theorem}

\begin{proof}
If $m\mid n$ with $m,n\in(P,2P]$ and $m\ne n$, then $n\ge 2m>2P$, contradicting $n\le 2P$.
The element $2(P+1)$ is the first proper multiple of an in-window element to exceed $2P$.
\end{proof}

\Formal{\lean{window\_self\_contained\_dvd}, \lean{self\_containment\_breaks\_just\_above}
(\lean{Rings.lean}); \lean{window\_self\_contained\_bound},
\lean{max\_self\_contained\_width} (\lean{SelfContained.lean}).}

Self-containment is not a threshold phenomenon that appears only for large bases; it is
universal, and a small window is not an exception but merely too narrow to contain such a
proper multiple. It is a structural property specific to the tightest window; \S\ref{sec:family}
shows that it is present exactly where the family is most constrained.

%======================================================================
\section{The generalized window family}
\label{sec:family}

\begin{definition}[Generalized window]
For $P_{\min},P_{\max}$ let $\mathrm{gW}(P_{\min},P_{\max}):=(P_{\max},P_{\min}P_{\max}]$
and say it \emph{has a void} if some $n\in\mathrm{gW}(P_{\min},P_{\max})$ is
$\mathcal{B}(P_{\max})$-void.
\end{definition}

\Formal{\lean{gwindow}, \lean{gwindowHasVoid} (\lean{Rings.lean}).}

\begin{proposition}[Bertrand is the case $P_{\min}=2$]\label{prop:two}
$\mathrm{gW}(2,P_{\max})=W(P_{\max})$, and $\mathrm{gW}(2,P_{\max})$ has a void iff
$W(P_{\max})$ contains a prime.
\end{proposition}

\Formal{\lean{gwindow\_two\_eq\_window}, \lean{gwindowHasVoid\_two\_iff}.}

The family is monotone in the width parameter: widening the window can only help.

\begin{theorem}[Monotonicity: tight $\Rightarrow$ wide]\label{thm:mono}
If $P_{\min}\le P_{\min}'$ then $\mathrm{gW}(P_{\min},P_{\max})\subseteq
\mathrm{gW}(P_{\min}',P_{\max})$; hence a void for the narrower window is a void for the
wider one. In particular every member of the family follows from the tightest member
$P_{\min}=2$.
\end{theorem}

\begin{proof}
$P_{\min}P_{\max}\le P_{\min}'P_{\max}$ gives the inclusion of half-open intervals with the
same left endpoint; a void position and its voidness are preserved.
\end{proof}

\Formal{\lean{gwindow\_subset}, \lean{gwindowHasVoid\_mono},
\lean{gwindowHasVoid\_of\_bertrand}.}

At the wide end the statement is free: once $P_{\min}P_{\max}$ is large enough, a prime in
the window is delivered by elementary means (a Euclid-type witness), with no counting.

\Formal{\lean{gwindowHasVoid\_of\_wide}.}

Two points deserve emphasis. The family is \emph{equivalent} to Bertrand, not strictly
stronger: by Theorem~\ref{thm:mono} the narrow case implies the wide cases, so the content of
the family is concentrated at $P_{\min}=2$. Moreover the structure is richest there:
self-containment (Theorem~\ref{thm:self}) holds \emph{only} for $P_{\min}=2$, since any wider
window already contains $2(P_{\max}+1)$, a proper multiple of an in-window element. The tightest
window thus carries the most structure --- precisely the case that is hardest to establish.

This is the place to state the scale invariance the family embodies. The invariant is not
a reflection of coordinates but a dilation: the object is the pair
$\big((X,2X],\ \{q\ \text{prime}:q<X\}\big)$, and under $X\mapsto\lambda X$ the window and
the base rescale \emph{together}, the ratio $2$ being preserved. Freezing the base at one
scale while sliding the window to another is not a configuration of the family at all ---
it mismatches the two dilations. Bertrand's $(n,2n]$ is simply the unit cell of this scale,
$P_{\min}=2$; motion along $\mathbb{N}$ is the renormalization $n\mapsto\lambda n$ with
window and base co-scaling. Proposition~\ref{prop:scale} is the invariance of the local
rule; the count that follows inherits it exactly (\S\ref{sec:atom},
\lean{coverFiber\_eq\_scaled}).

%======================================================================
\section{Truncated arithmetic: the reduced-base sieve}
\label{sec:trunc}

The family of \S\ref{sec:family} is produced by \emph{truncated arithmetic}: instead of
the full base one sieves with the reduced base of primes that can actually reach the
window. By Lemma~\ref{lem:lpf} the only rings that cover a composite $n\le 2P$ are those
with $p\le\sqrt{2P}$; equivalently $p^2\le 2P$. This selects a truncated primorial.

\begin{definition}[Truncated primorial]
$\displaystyle M_{\mathrm{tr}}(P):=\prod_{\substack{p<P\ \text{prime}\\ p^2\le 2P}}p.$
\end{definition}

\Formal{\lean{truncPrimorial}; general version \lean{gtruncPrimorial}.}

\begin{theorem}[Void through the truncated base]\label{thm:trunc}
For $2<P$ and $n\in W(P)$, $n$ is $\mathcal{B}(P)$-void iff $\gcd(n,M_{\mathrm{tr}}(P))=1$.
Consequently $W(P)$ has a void iff $(P,2P]$ contains an integer coprime to
$M_{\mathrm{tr}}(P)$.
\end{theorem}

\begin{proof}
Coprimality to the full base clearly implies coprimality to the truncated base. Conversely,
if $n$ is coprime to $M_{\mathrm{tr}}(P)$ but divisible by some prime $p<P$, then $n$ is
composite, so by Lemma~\ref{lem:lpf} $\operatorname{minFac}(n)^2\le 2P$, and
$\operatorname{minFac}(n)$ divides $M_{\mathrm{tr}}(P)$ --- contradiction.
\end{proof}

\Formal{\lean{void\_iff\_coprime\_trunc}, \lean{windowHasVoid\_iff\_trunc\_coprime},
\lean{survivor\_of\_trunc\_coprime}; general \lean{gvoid\_iff\_coprime\_trunc},
\lean{gwindowHasVoid\_iff\_trunc\_coprime}.}

The truncation principle makes the reach explicit at every level and exposes the
self-similarity. Removing only the prime $2$ leaves reach $3n$; removing $2$ and $3$ leaves
$5n$; in general the smallest omitted prime dictates the reach. The buildability map is
therefore layered: the class of positions with least factor $p$ is buildable up to
$p\,P_{\max}$, and the same phenomenon recurs at each level with Bertrand as the base case
$P_{\min}=2$.

\Formal{Reduced-base sieve and its non-closure:
\lean{exists\_uncovered\_of\_card\_lt}, \lean{truncated\_not\_sieve\_closed},
\lean{card\_truncated\_window} (\lean{Truncated.lean}).}

%======================================================================
\section{Reduction to a single window-count atom}
\label{sec:atom}

We now count. Let $\mathrm{cov}(P)$ be the number of composite (covered) positions in
$W(P)$; since $|W(P)|=P$, non-emptiness is a strict inequality.

\begin{definition}[Covered count]
$\mathrm{cov}(P):=\#\{n\in W(P): n\ \text{is not}\ \mathcal{B}(P)\text{-void}\}.$
\end{definition}

\begin{theorem}[Reduction to an atom]\label{thm:atom}
$W(P)$ has a void iff $\mathrm{cov}(P)<P$. Hence Bertrand at scale $P$ is equivalent to the
single positivity statement $P-\mathrm{cov}(P)>0$.
\end{theorem}

\Formal{\lean{coveredCount}, \lean{gps\_window\_card},
\lean{windowHasVoid\_iff\_coveredCount\_lt}, \lean{coveredCount\_lt\_imp\_nonempty}.}

A structural preliminary records that the largest ring $P_{\max}$ exerts no force inside
its own window (its first multiple $2P_{\max}$ is the lone edge point), and that the
previous base carries a weight $w=P_{\max}\cdot\varphi(M)/M\ge1$, proved once and for all,
independently of scale. This weight is structural content about the setting, not a
hypothesis either regime's closure consumes below: the sparse regime closes by a union
bound (Theorem~\ref{thm:jac} and the paragraph preceding it), and the dense regime closes
on the central binomial coefficient (\S\ref{sec:erdos}), neither appealing to $w\ge1$. It is
recorded here because it is the reduction's own by-product, not because the atom's proof
term invokes it.

\Formal{\lean{zero\_effective\_force} (\lean{ZeroForce.lean}); weight
$w\ge1$: \lean{WeightGeOne}, \lean{weight\_base}, \lean{weight\_step},
\lean{structural\_weight\_ge\_one}, \lean{prev\_base\_cannot\_cover\_window}
(\lean{Weight.lean}); bridge \lean{weight\_density\_bridge}; anchor idle
\lean{anchor\_idle\_in\_own\_window}. This chain is proved and wired through
\lean{structural\_sieve\_survivor}'s signature, but \lean{dense\_sieve\_survivor} --- the
lemma that actually closes the dense regime --- takes the resulting density fact as an
unused parameter (\lean{\_hdensity}) and proves the window nonempty directly from
\lean{binomial\_contradiction} instead.}

The covered count decomposes disjointly by least prime factor, and each fiber is, by an
exact bijection, the \emph{same} coprimality count one scale down --- the recursion engine
promised by scale invariance.

\begin{theorem}[Disjoint minFac decomposition and self-similar fibers]\label{thm:fiber}
$\displaystyle \mathrm{cov}(P)=\sum_{p}\, A_p(P)$, where $A_p(P)$ counts the $n\in W(P)$
with $\operatorname{minFac}(n)=p$, the sum ranging over primes with $p^2\le 2P$. Moreover
\[
A_p(P)=\#\Big\{m\in\big(\tfrac{P}{p},\tfrac{2P}{p}\big]:
\gcd\big(m,\textstyle\prod_{r<p}r\big)=1\Big\},
\]
the same window-coprimality problem at scale $P/p$.
\end{theorem}

\Formal{\lean{coveredCount\_eq\_sum\_minFac\_fiber},
\lean{coveredCount\_eq\_sum\_truncated}, \lean{coverFiber\_eq\_scaled};
the dominant ring $p=2$ is split off exactly in \lean{coveredCount\_split\_two},
\lean{windowHasVoid\_iff\_oddSum\_lt}.}

Two exact handles on $\mathrm{cov}(P)$ follow. A union bound with multiplicity gives
$\mathrm{cov}(P)\le\sigma(P):=\sum_{p}(\text{hits of }p)$; the gap between the distinct
count $\mathrm{cov}(P)$ and $\sigma(P)$ is exactly the overlap (redundancy) between rings,
so $\sigma(P)>P$ does not obstruct $\mathrm{cov}(P)<P$. The full, signed identity is
inclusion--exclusion (Legendre).

\begin{theorem}[Interference identity]\label{thm:legendre}
For a finite set $S$ of primes and $a\le b$,
\[
\#\Big\{n\in(a,b]:\gcd\big(n,\textstyle\prod_{q\in S}q\big)=1\Big\}
=\sum_{T\subseteq S}(-1)^{|T|}
\Big(\big\lfloor\tfrac{b}{\prod_{q\in T}q}\big\rfloor
-\big\lfloor\tfrac{a}{\prod_{q\in T}q}\big\rfloor\Big).
\]
The order-by-order layers are its partial sums; the whole is an exact identity.
\end{theorem}

\Formal{\lean{coveredCount\_le\_sigma}, \lean{interference\_step},
\lean{interference\_formula} (\lean{Rings.lean}).}

Entire regimes of the atom close with no analysis. In the \emph{sparse} regime a single
union bound produces a survivor. For small anchors a sufficient condition settles the
window: only the \emph{active-covering rings} $p^2\le 2P$ can cover a composite in $W(P)$
(Lemma~\ref{lem:lpf}), so, writing $M_{\mathrm{tr}}(P)=\prod_{p^2\le 2P}p$, if every run of
$P$ consecutive integers meets a number coprime to $M_{\mathrm{tr}}(P)$ --- i.e.\ the
coprime-gap $g\big(M_{\mathrm{tr}}(P)\big)\le P$ --- then $W(P)$ contains a coprime, hence a
void.

\begin{theorem}[Small-anchor closure]\label{thm:jac}
If $g\big(M_{\mathrm{tr}}(P)\big)\le P$ then $W(P)$ has a void. Concretely $g(30)=6$,
$g(210)=10$, $g(2310)=14$ close all anchors $P\le 83$ unconditionally.
\end{theorem}

A caveat on this bound. Here $g\big(M_{\mathrm{tr}}(P)\big)$ is the global coprime-gap of a
\emph{fixed} modulus --- the product of the active-covering rings, a subcollection that
coincides with a smaller, foreign base, not the complete base of the object. It only
\emph{upper-bounds} the object's own quantity, the local run of covered positions inside
$W(P)$; the inequality $g\le P$ is therefore a loose sufficient proxy, not the governing
structure of the object, and it ceases to be available at large $P$, where the residual
dense regime is settled by the object below, not by any gap function.

\Formal{\lean{windowHasVoid\_of\_jacobsthal}, \lean{windowHasVoid\_of\_trunc\_gap},
\lean{jacobsthal\_thirty}, \lean{jacobsthal\_210}, \lean{jacobsthal\_2310}, and
\lean{windowHasVoid\_3} \dots \lean{windowHasVoid\_83} (\lean{Rings.lean}).}

This closure is independent of the assembled proof, not a stage feeding into it:
\lean{structural\_sieve\_survivor} never branches on $P\le83$, and dispatches every
prime $P$ through the sparse/dense split of \S\ref{sec:atom}, the dense side closed
uniformly by \lean{binomial\_contradiction} at the single threshold $n=512$
(\S\ref{sec:erdos}). Theorem~\ref{thm:jac} thus supplies a second, self-contained
certificate for these finitely many scales, disjoint in its derivation from
Theorem~\ref{thm:window}; it is retained to exhibit a regime that closes with no
counting at all, not because the main theorem's proof term invokes it.

What remains is the \emph{dense} regime at large $P$. Here the structural mechanism has a
definite reach and a definite limit, and it is worth stating both precisely. At every
\emph{concrete} scale $P$ the construction decides the window outright: the sieve is
finite and deterministic, so each finite step is verified and no counterexample can arise
--- the statement \enquote{$W(P)$ has a void} is decidable at every fixed $P$. What the
generative process cannot deliver is the \emph{uniform} statement over all scales at once:
to reach \enquote{for every $P$} by generation would mean carrying the construction to
infinity, an infinite process that no finite computation realizes. The passage from
\enquote{true at every finite step} to \enquote{true in the infinite} is therefore not
computational but \emph{existential}. We do not claim a generative closure of this
existential step, and none is known. It is precisely this closure of the infinite tail ---
the part beyond any finite computation --- that the central binomial coefficient performs, and
the next section explains why it is the appropriate object for this step.

\Formal{Regime dispatch: \lean{structural\_sieve\_survivor}; the residual dense case
\lean{dense\_sieve\_survivor} (\lean{GPS\_StateMachine.lean}).}

%======================================================================
\section{The central binomial coefficient and the window's prime content}
\label{sec:newton}

The dense regime is closed existentially, and the central binomial coefficient is the
instrument of that closure. It is invoked for one purpose only: to settle the infinite tail
that no finite generation can reach --- a uniform argument in place of an unrealizable infinite
process, not a repair of the finite mechanism, which needs none, since that mechanism already
decides every concrete scale. The coefficient need not be present for any fixed window; it is
present because the proof must quantify over all windows at once, and it supplies closure over
that range.

That $\binom{2n}{n}$ is the appropriate object here is itself structural. It is an entry of
Newton's binomial expansion: $4^n=(1+1)^{2n}$ is the sum of the $2n+1$ entries of the $2n$-th
Pascal row, of which the central one is largest, and the window's new sources are recorded in
its prime factorization by a purity law internal to the construction.

\begin{theorem}[$S1$ purity law]\label{thm:s1}
A new source $q\in(n,2n]$ divides $\binom{2n}{n}$ exactly once: its first multiple $2q$
already exceeds $2n$, so $q$ contributes no further factor to $(2n)!$, while $q>n$ keeps it out
of $(n!)^2$.
Consequently the product of all new sources of the window divides $\binom{2n}{n}$:
\[
\prod_{\substack{q\in(n,2n]\\ q\ \text{prime}}} q \ \Big|\ \binom{2n}{n},
\qquad\text{hence}\qquad
\prod_{\substack{q\in(n,2n]\\ q\ \text{prime}}} q \ \le\ \binom{2n}{n}.
\]
\end{theorem}

\Formal{\lean{window\_primes\_prod\_dvd\_centralBinom} (\lean{Rings.lean}),
\lean{window\_primes\_prod\_dvd\_centralBinom'},
\lean{window\_primes\_prod\_le\_centralBinom} (\lean{Newton.lean}); the
divisibility rests on \lean{prod\_primes\_dvd} and
\lean{Nat.Prime.dvd\_choose\_add}. The purity is the same self-containment property as
\lean{window\_self\_contained\_dvd} / \lean{window\_composite\_smooth} --- the first proper
multiple of an in-window source already exceeds $2P_{\max}$ --- with the same role of $2$ and
the same extremal self-contained case $P_{\min}=2$.}

The choice of object is thus forced by the structure, not imported from history: the window's
new sources are recorded in the prime factorization of $\binom{2n}{n}$. Its role in the proof
is existential --- it closes what cannot be generated --- while its identity is structural.
Two bounds close on it. The lower bound is pure combinatorics of the Pascal row.

\begin{theorem}[Lower bound]\label{thm:lower}
$4^n\le(2n+1)\binom{2n}{n}$, proved directly from
$4^n=\sum_{i=0}^{2n}\binom{2n}{i}$ with each term $\le\binom{2n}{n}$.
\end{theorem}

\Formal{\lean{four\_pow\_le\_newton} (\lean{Newton.lean}).}

The upper bound reads the same factorization in the contrapositive: an empty window would
force $\binom{2n}{n}$ to be built entirely from old sources.

\begin{theorem}[Empty window $\Rightarrow$ old sources only]\label{thm:upper}
If $(n,2n]$ contains no prime, then every prime factor of $\binom{2n}{n}$ is $\le n$.
\end{theorem}

\begin{proof}
Any prime factor of $\binom{2n}{n}$ divides $(2n)!$, hence is $\le 2n$; absent new sources
in $(n,2n]$ it is $\le n$.
\end{proof}

\Formal{\lean{centralBinom\_prime\_factor\_le} (\lean{Newton.lean}).}

The purity law and the two bounds above (Theorems~\ref{thm:s1}--\ref{thm:upper}) are proved
in \lean{Rings.lean} and \lean{Newton.lean} and were the construction's first route to
closure; they remain true and are kept in the development. The main theorem's proof term,
however, closes the dense regime by a later, independent route (\S\ref{sec:erdos}): a
factorization upper bound reproved from scratch in \lean{BinomialBound.lean}, combined with
Mathlib's own central-binomial lower bound, rather than through \lean{Newton.lean}'s
in-project lower bound or \lean{Rings.lean}'s purity-law upper bound. Theorems
\ref{thm:s1}--\ref{thm:upper} are retained as an alternative derivation of the same closing
idea, not because \lean{binomial\_contradiction} (\S\ref{sec:erdos}) invokes them.

%======================================================================
\section{Closing the extremal window: the self-contained certificate}
\label{sec:erdos}

Only the extremal window $P_{\min}=2$ needs closing, and only in the dense regime. We
split at $n=512$.

For $2<n<512$ a kernel-checked decision procedure exhibits a prime in $(n,2n]$ directly: the
range is split into finitely many blocks, and on each a witness prime is exhibited and
verified by direct reduction inside Lean's trusted kernel --- no external lemma, and no
appeal to compiled evaluation.

\Formal{\lean{small\_window\_oracle}, by \lean{decide} (the range split via
\lean{interval\_cases} into blocks the kernel evaluator can handle in one call).}

For $n\ge512$ the closure uses a separate pair of bounds from \S\ref{sec:newton}'s
purity-law route: Mathlib's own central-binomial lower bound, and an old-source upper
bound on $\binom{2n}{n}$ reproved independently in \lean{BinomialBound.lean} (at most
$(2n)^{\sqrt{2n}}\,4^{2n/3}$, the same shape as Theorem~\ref{thm:upper} but not that
theorem). Suppose $(n,2n]$ has no prime; combining this lower bound, this upper bound, and
the size threshold below yields $4^n<4^n$, a contradiction.

\begin{theorem}[Quantitative kernel]\label{thm:erdos}
For $2<n$, if no prime lies in $(n,2n]$ then a contradiction follows. Hence every window
$(n,2n]$ with $n>1$ contains a prime.
\end{theorem}

\begin{proof}[Proof sketch]
For $n<512$ use the oracle. For $n\ge512$, chain the lower bound
$4^n<n\binom{2n}{n}$ (\lean{Nat.four\_pow\_lt\_mul\_centralBinom}), the old-source upper
bound $\binom{2n}{n}\le(2n)^{\sqrt{2n}}4^{2n/3}$ (\lean{window\_centralBinom\_le}, reproved
in-project from Legendre/Kummer and primorial primitives), and the prime-free size
comparison $n\,(2n)^{\sqrt{2n}}4^{2n/3}\le4^n$ (\lean{threshold\_inequality}) to reach
$4^n<4^n$.
\end{proof}

\Formal{\lean{binomial\_contradiction} (\lean{BinomialCertificate.lean}). The upper bound
(\lean{window\_centralBinom\_le}, \lean{BinomialBound.lean}) is reproved from
\lean{Choose.Factorization} and \lean{Primorial}; the size threshold
(\lean{threshold\_inequality}, \lean{Threshold.lean}) from real convexity. Neither imports
\lean{Mathlib.NumberTheory.Bertrand}, so the main theorem's import closure is free of it.
The certificate is modular (\lean{WindowCertificate}, \lean{Certificate.lean}): any estimate
of the same strength --- a Chebyshev-type bound, or Nair's $\operatorname{lcm}(1,\dots,2n)$
argument \cite{nair1982} --- plugs into the identical slot; a second instance,
\lean{erdos\_certificate}, does so through Mathlib's inequalities.}

Non-circularity is by inspection: Mathlib's \lean{Nat.bertrand} and
\lean{exists\_prime\_lt\_and\_le\_two\_mul} are never used, and
\lean{Mathlib.NumberTheory.Bertrand} is not in the import closure of the main theorem; it
appears only in \lean{Erdos.lean}, which provides the instance-A comparison and is not
imported by the main theorem.

%======================================================================
\section{Main theorem and recovery of the classical case}
\label{sec:main}

Assembling the reduction (\S\ref{sec:atom}), the factorization identity (\S\ref{sec:newton})
and the certificate (\S\ref{sec:erdos}) closes the atom at every scale, hence produces a prime
in every window.

\begin{theorem}[Window non-emptiness]\label{thm:window}
For every prime $P$ with $2<P$, the window $(P,2P]$ contains a prime.
\end{theorem}

\Formal{\lean{structural\_sieve\_survivor} $\to$ \lean{gps\_window\_nonempty} $\to$
\lean{prime\_in\_window} (\lean{GPS\_StateMachine.lean}).}

\begin{theorem}[Generalized family]\label{thm:gen}
For every prime $P_{\max}$ and every $P_{\min}\ge2$, the window
$(P_{\max},P_{\min}P_{\max}]$ contains a prime.
\end{theorem}

\begin{proof}
By Theorem~\ref{thm:mono} it suffices to treat $P_{\min}=2$, which is
Theorem~\ref{thm:window}.
\end{proof}

\Formal{\lean{gwindowHasVoid\_of\_bertrand}.}

\begin{corollary}[Classical Bertrand--Chebyshev]\label{cor:classical}
For every integer $N>1$ there is a prime $p$ with $N<p\le2N$.
\end{corollary}

\begin{proof}
Let $P$ be the largest prime $\le N$. If $P=2$ then $N=2$ and $p=3$ works. Otherwise
$2<P$; Theorem~\ref{thm:window} gives a prime $q\in(P,2P]$, and maximality of $P$ forces
$q>N$, while $P\le N$ gives $q\le2N$.
\end{proof}

\Formal{\lean{exists\_largest\_prime\_le}, \lean{bertrand\_chebyshev} (\lean{Main.lean});
the base-consistent structural form is \lean{structural\_bertrand\_chebyshev}.}

The constant $2$ is not arbitrary: it is the smallest prime, hence the smallest admissible
$P_{\min}$, hence the buildability edge (\S\ref{sec:edge}). Bertrand's window is thus the
extremal member of the family --- the case $P_{\min}=2$, on which all wider members depend and
at which the self-containment structure is present.

%======================================================================
\section{Remarks and provenance}
\label{sec:provenance}

\paragraph{Comparison with prior work.}
The construction's closest relatives in the literature are Jacobsthal's function $g(n)$,
the largest gap between consecutive integers coprime to the primorial of the first $n$
primes \cite{jacobsthal1961}, later sharpened by Iwaniec \cite{iwaniec1978}, and the
generalized Bertrand-type window results of Das and Paul, who prove existence of a prime
in $[n,kn]$ for general $k$ by direct counting and sieve estimates \cite{daspaul2019}. Both
are thematically adjacent: Jacobsthal's function bounds gaps in exactly the reduced-base
setting this paper's truncated arithmetic (\S\ref{sec:trunc}) sieves with, and Das--Paul's
result generalizes the same window statement this paper's family (\S\ref{sec:family})
assembles. Neither, to the author's knowledge, derives the window and its endpoint from the
buildability edge of a prime base, or reduces the postulate to the single exact-count atom
of \S\ref{sec:atom}; a search of the literature on Jacobsthal's function, on elementary and
generalized proofs of Bertrand's postulate, and on primorial-coprimality sieve arguments did
not surface a prior source assembling these three components in this form. This search is
necessarily bounded --- it covers indexed literature (arXiv, standard references on
Jacobsthal's function and Bertrand's postulate) and cannot rule out unindexed material such
as course notes or unpublished manuscripts --- so the claim above is offered as a
literature-grounded "to the author's knowledge," not as an assertion of absolute priority.

\paragraph{Machine verification.}
The development is checked in Lean~4 with Mathlib (toolchain pinned in the repository);
there is no \texttt{sorry}, and the main theorem depends on no axiom beyond the three
standard ones (\texttt{propext}, \texttt{Classical.choice}, \texttt{Quot.sound}), verified
directly by \lean{\#print axioms bertrand\_chebyshev} --- in particular
\lean{Lean.ofReduceBool} (the axiom \lean{native\_decide} introduces) does not appear
anywhere in its import closure. The small-window oracle for $n<512$, the base anchors
$P\le83$, and the Jacobsthal gap values are all discharged by \lean{decide}, fully reduced
and checked by Lean's own kernel. The two widest of these --- the $n<512$ search and the
$g(2310)$ gap check --- exceed the kernel's evaluator as a single call and are each split
into blocks small enough for it to handle, rather than falling back to compiled evaluation.
(A separate, Mathlib-based comparison instance in \lean{Erdos.lean} does use
\lean{native\_decide}; it lies outside the main theorem's import closure and does not affect
this.) A non-circularity audit (a search over uses) confirms that Mathlib's own Bertrand
theorem is not invoked.

\paragraph{Provenance.}
The development is neither independent of Erd\H{o}s nor Erd\H{o}s in disguise. The
structural reduction is independent: it derives the window and its edge from the base
(\S\ref{sec:reach}--\S\ref{sec:edge}), assembles the scale-invariant family
(\S\ref{sec:family}--\S\ref{sec:trunc}), and localizes the theorem to one exact-count atom
whose small, sparse and gap regimes close with no counting (\S\ref{sec:atom}). The atom is
settled on the central binomial coefficient, whose prime factorization records the window's
content via the $S1$ purity law (\S\ref{sec:newton}) and would be reached here regardless of
the historical record. Erd\H{o}s reached the same object in 1932 by a different route: a
convergence on one mathematical object, carrying historical priority of use but no exclusivity
over it. In the Lean development the two bounds are reproved in-project, so the main theorem's
import closure contains no Mathlib Bertrand theorem; the certificate is modular
(\S\ref{sec:erdos}), a second, Mathlib-based instance being retained, outside that import
closure, for comparison.

\paragraph{Relation to classical proofs.}
The successive treatments of the postulate answered different questions and, accordingly,
studied different objects, while converging on one theorem. Bertrand \cite{bertrand1845} posed
the question of existence, read from tables of the integers; Chebyshev \cite{chebyshev1852}
answered it through the analytic size of $\vartheta$ and $\psi$; Erd\H{o}s \cite{erdos1932} and
Nair \cite{nair1982} through the elementary size of $\binom{2n}{n}$ and of
$\operatorname{lcm}(1,\dots,2n)$. The present treatment asks a fourth question --- the
deterministic reach of a generating base --- and takes the base, not the window, as its primary
object; the window and its endpoint are then derived, and the existence of a prime is a
corollary of the reach rather than the premise. At the single quantitative point all four are
of the same strength, and any of the elementary certificates plugs into the same slot; what the
structural layer adds is the reduction that makes one such certificate sufficient, together with
the derivation of the window itself from the base.

\paragraph{Further directions.}
The finite mechanism decides every concrete scale; what a certificate supplies is the
existential closure of the infinite tail. Whether that closure can be effected without a
size inequality is open; every method known to the author is quantitative, of the same
Chebyshev strength.

%======================================================================
\section*{Competing Interests}
The author declares no competing interests.

\section*{Funding}
No funding was received for this work.

\printbibliography

\end{document}
